\color{unimain}{\subsection{Автоматическое ускорение при включении выключателя}}
\color{black}

\begin{enumerate}

\item
В устройстве реализована логика автоматического ускорения (АУ) ДЗ, ТНЗНП и МТЗ при включении выключателя.

\item
Уставки АУ представлены в таблице \ref{app_set:au}. На рисунке \ref{pict_au:AU} приведен алгоритм работы АУ. 


\vspace{2mm}
\begin{figure}[H]
\centering
\includegraphics[width=0.8\textwidth,height=0.8\textheight,keepaspectratio]{\fbpath/07.02 РЗАТ/041.1351 АУ/img/AU.pdf}
\captionof{figure}{Функциональная схема логики АУ}\label{pict_au:AU}
\end{figure}

\item
Логика автоматического ускорения, в зависимости от первичной схемы подключения АТ к РУ, вводится:
\begin{itemize}
\item АУ в сторону АТ;
\item АУ в сторону сети;
\item Ненаправленное АУ.
\end{itemize}

\item
Ввод в работу автоматического ускорения ступеней, направленных в сторону АТ, осуществляется программным ключом \textbf{«АУ в АТ»}. Выдержка времени срабатывания ускоряемой ступени ДЗ в АТ регулируется уставкой \textbf{«Тср АУ ДЗ в АТ»}, ступени ТНЗНП в АТ -- \textbf{«Тср АУ ТНЗНП в АТ»}, ступени МТЗ в АТ -- \textbf{«Тср АУ МТЗ в АТ»}. АУ в сторону АТ с контроль отсутствия напряжения задается программным переключателем \textbf{«КОН АУ в АТ»}. Условия ввода АУ в АТ с КОН -- ускорение по факту включения выключателя своей стороны и отсутствию напряжения на АТ (цепи напряжения от ТН НН) или при отключенном состоянии выключателей смежной стороны. 

\item
Условия разрешения КОН для АУ в сторону АТ задаются программным переключателем \textbf{«Режим КОН АТ»}:
\begin{itemize}
\item \textit{«ТН НН»} -- Разрешение АУ в сторону АТ при отсутствии напряжения ТН НН, сигнал \textit{«Ср. КОН НН»};
\item \textit{«РПО см»} -- Разрешение АУ в сторону АТ при отключенном состоянии смежной стороны, сигнал \textit{«См. ст. отключена»};
\item \textit{«ТН НН или РПО см»} -- Разрешение АУ в сторону АТ при отсутствии напряжения ТН НН или при отключенном состоянии смежной стороны.
\end{itemize}

\item
Ввод в работу автоматического ускорения ступеней, направленных в сторону сети, осуществляется программным ключом \textbf{«АУ в сеть»}. Выдержка времени срабатывания ускоряемой ступени ДЗ в сеть регулируется уставкой \textbf{«Тср АУ ДЗ в сеть»}, ступени ТНЗНП в сеть -- \textbf{«Тср АУ ТНЗНП в сеть»}, ступени МТЗ в сеть -- \textbf{«Тср АУ МТЗ в сеть»}. АУ в сторону сети с контроль отсутствия напряжения задается программным переключателем \textbf{«КОН АУ в сеть»}. Условия ввода АУ в сеть с КОН -- ускорение по факту включения выключателя своей стороны и отсутствию напряжения на шинах (цепи напряжения от ТН СШ).

\item
Ввод в работу автоматического ускорения ненаправленных ступеней, осуществляется программным ключом \textbf{«Ненапр АУ»}. Выдержка времени срабатывания ускоряемой ненаправленной ступени ДЗ в сеть регулируется уставкой \textbf{«Тср АУ ДЗ ненапр»}, ненаправленной ступени ТНЗНП в сеть -- \textbf{«Тср АУ ТНЗНП ненапр»}, ненаправленной ступени МТЗ в сеть -- \textbf{«Тср АУ МТЗ ненапр»}. АУ ненаправленных ступеней с
контроль отсутствия напряжения задается программным переключателем \textbf{«КОН АУ ненапр»}. Условия ввода АУ от ненаправленных ступени c КОН -- ускорение по факту включения выключателя своей стороны и отсутствии напряжения на СШ или на АТ.

\item
Время ввода автоматического ускорения по факту включения выключателя задается уставкой \textbf{«Тввод АУ»}.

\item
Оперативный ввод/вывод АУ производится с помощью БУ (блоков управления функциями) \textit{«Вывод АУ»}.

\end{enumerate}